\documentclass[11pt]{article}
\usepackage{acl}
\usepackage{times}
\usepackage{latexsym}
\usepackage[T1]{fontenc}
\usepackage[utf8]{inputenc}
\usepackage{microtype}
\usepackage{inconsolata}
\usepackage{graphicx}
\usepackage{float}
\usepackage{caption}
\usepackage{booktabs}


\captionsetup{skip=4pt}

\setlength{\textfloatsep}{8pt}   % space above/below floats
\setlength{\floatsep}{6pt}       % space between floats
\setlength{\intextsep}{6pt}      % space for in-text floats

\title{Evaluating DesignBench Baselines for Automated UI Code Repair}

\author{Alice Le \\
  Carnegie Mellon University \\
  \texttt{anhle@andrew.cmu.edu} \\\And
  Isaac Au \\
  Carnegie Mellon University \\
  \texttt{yau@andrew.cmu.edu} \\}

\begin{document}
\maketitle

\begin{abstract}
We reproduce Qwen2.5-VL (72B and 7B) on the DesignBench UI code repair task and conduct a systematic error analysis to inform our proposed approach of repurposing GUI grounding models as visual critics. Our 72B reproduction closely matches reported results, with Angular CMLS and CMCS within 5\%. The 7B model exhibits severe compile failures on framework-based projects (75\% on React, 43\% on Angular) and low Issue Accuracy (27\%), motivating the 72B model as our generation backbone. Error analysis reveals that defect identification and code localization are compounding bottlenecks, and that spatial defects (overflow, occlusion) are substantially harder than property-level defects (color, contrast). These findings refine our proposal: we use 72B for generation while targeting the identification and localization gaps with GUI grounding augmentation.
\end{abstract}


%% ─────────────────────────────────────────────────────────────
\section{Baseline Reproduction}
\label{sec:baseline}
   % Rubric: existing baseline? describe it, SOTA justification, results + differences

\subsection{Task and Baseline}
% - DesignBench Design Repair task description (input: screenshot + buggy code + issue list → repaired code)
% - Metrics: CSR, CMLS, CMCS
% - Baseline choice: Qwen2.5-VL (72B + 7B)
% - SOTA justification: GPT-4o is SOTA but proprietary; Qwen is top open-weight model, two sizes enable scale analysis

We reproduce the \textit{Design Repair} subtask from DesignBench \cite{designbench2025}, a benchmark for evaluating multimodal LLMs on automated front-end engineering. Given a rendered screenshot of a broken UI, its buggy source code, and a list of display issues, a model must output repaired code that resolves the visual defects. DesignBench covers four front-end frameworks, namely React, Vue, Angular, and vanilla HTML/CSS. There are six defect categories: misalignment, overflow, occlusion, crowding, color/contrast, and other. The repair split contains 111 samples distributed across the four frameworks (28/27/28/28).
 
\paragraph{Metrics.} DesignBench evaluates repair quality along five dimensions CSR, IssAcc, CMLS, CMCS, CLIP as defined in Table~\ref{tab:metrics}.
 
\begin{table}[t]
\centering
\small
\begin{tabular}{lll}
\toprule
\textbf{Metric} & \textbf{Measures} & \textbf{Stage} \\
\midrule
CSR     & Code compiles without errors       & Generation \\
IssAcc  & Correct defect-type labels         & Identification \\
CMLS    & Jaccard sim.\ of AST edits         & Localization \\
CMCS    & CMLS $\times$ CodeBLEU of edits    & Loc.\ + Content \\
CLIP    & Visual sim.\ of rendered output    & Overall \\
\bottomrule
\end{tabular}
\caption{DesignBench metrics mapped to repair pipeline stages.}
\label{tab:metrics}
\end{table}


\paragraph{SOTA Justification \& Baseline choice.}
DesignBench provides the direct existing baseline for our task. The DesignBench paper evaluates nine MLLMs on the repair task, spanning proprietary models (GPT-4o, Claude-3.7-Sonnet, Gemini-2.0-Flash) and open-weight models (Qwen2.5-VL 72B/7B, Llama-3.2 90B/11B, Pixtral 124B/12B). Among models evaluated in the benchmark, GPT-4o and
Claude-3.7-Sonnet achieve the strongest overall performance. Since the benchmark's release, newer proprietary models (Claude Opus 4.6, GPT-5.4) have become available and likely outperform the evaluated models, but have not been tested on DesignBench. All proprietary models---both those in the original evaluation and their successors---share the same reproducibility limitation that API outputs are non-deterministic and underlying checkpoints change without notice.

Among open-weight models evaluated on DesignBench, Qwen2.5-VL-72B is the strongest performer on the repair task. It is available via the Dashscope API---the same inference endpoint used by the original authors, making it the most reproducible strong baseline available. We additionally evaluate Qwen2.5-VL-7B for two reasons. First, it enables a within-family scale analysis. Second, it is the base model for JEDI-7B
\cite{xie2025jedi}, the GUI grounding model central to our proposed approach; characterizing the 7B model's failure modes is essential
groundwork for understanding where grounding augmentation can provide the most benefit.

We use the same prompt template as the original paper (\texttt{run\_repair.py} $\to$ \texttt{repair\_prompt.py}), running in \texttt{both} mode (screenshot + code) throughout. 
 
\subsection{Reproduction Results}
 % - Tables: 72B and 7B results, ours vs paper
% - Discussion of differences: broadly consistent, CSR gaps are detection artifacts (regex version mismatch), small N per framework
% - Note: CLIP visual similarity confirms comparable outputs
% - Remaining unexplained variance: possible Dashscope API checkpoint drift, but unverifiable
\begin{table*}[t]
\centering
\small
\setlength{\tabcolsep}{4pt}
\begin{tabular}{lrrrrrr}
\toprule
\textbf{FW} & \textbf{N} & \textbf{CSR} & \textbf{CMLS} & \textbf{CMCS} & \textbf{IssAcc} & \textbf{CLIP} \\
 & & \multicolumn{3}{c}{\footnotesize (paper / ours)} & \multicolumn{2}{c}{\footnotesize (ours)} \\
\midrule
React   & 28 & 0.929 / 0.857 & 0.634 / 0.339 & 0.542 / 0.230 & 0.395 & 0.771 \\
Vue     & 27 & 1.000 / 1.000 & 0.509 / 0.213 & 0.392 / 0.143 & 0.213 & 0.796 \\
Angular & 28 & 0.929 / 0.964 & 0.676 / 0.631 & 0.571 / 0.556 & 0.379 & 0.821 \\
Vanilla & 28 & ---\phantom{0} / 1.000 & 0.609 / 0.532 & 0.586 / 0.510 & 0.369 & 0.791 \\
\midrule
Overall & 111 & 0.952 / 0.955 & 0.607 / 0.430 & 0.523 / 0.360 & 0.340 & 0.795 \\
\bottomrule
\end{tabular}
\caption{Qwen2.5-VL-72B repair results. CSR, CMLS, and CMCS are shown as paper\,/\,ours; IssAcc and CLIP are ours only (the original paper reports IssAcc aggregated differently and does not report per-framework CLIP for repair).}
\label{tab:72b}
\end{table*}

\begin{table*}[t]
\centering
\small
\setlength{\tabcolsep}{4pt}
\begin{tabular}{lrrrrrr}
\toprule
\textbf{FW} & \textbf{N} & \textbf{CSR} & \textbf{CMLS} & \textbf{CMCS} & \textbf{IssAcc} & \textbf{CLIP} \\
 & & \multicolumn{3}{c}{\footnotesize (paper / ours)} & \multicolumn{2}{c}{\footnotesize (ours)} \\
\midrule
React   & 28 & 0.286 / 0.250 & 0.362 / 0.182 & 0.257 / 0.139 & 0.345 & 0.632 \\
Vue     & 27 & 0.111 / 0.926 & 0.245 / 0.237 & 0.180 / 0.179 & 0.210 & 0.785 \\
Angular & 28 & 0.036 / 0.571 & 0.641 / 0.304 & 0.577 / 0.206 & 0.173 & 0.486 \\
Vanilla & 28 & ---\phantom{0} / 1.000 & 0.499 / 0.422 & 0.478 / 0.394 & 0.345 & 0.796 \\
\midrule
Overall & 111 & ---\phantom{0} / 0.687 & 0.437 / 0.287 & 0.373 / 0.230 & 0.269 & 0.674 \\
\bottomrule
\end{tabular}
\caption{Qwen2.5-VL-7B repair results. Format matches Table~\ref{tab:72b}.}
\label{tab:7b}
\end{table*}

Tables~\ref{tab:72b} and~\ref{tab:7b} report our results alongside the paper's values for 72B and 7B respectively.
 

 

The 72B results are broadly consistent with the paper. Angular CMLS
(0.676 vs.\ 0.631) and CMCS (0.571 vs.\ 0.556) are within 5 percentage
points. Vue 7B CMCS is nearly identical (0.180 vs.\ 0.179). CLIP visual
similarity scores are comparable across all conditions, confirming that
model outputs are qualitatively similar even where code-level metrics
diverge.

The 7B CSR shows the largest discrepancies: Vue (paper 0.111 vs.\ ours
0.926) and Angular (0.036 vs.\ 0.571). We verified that Node versions,
framework versions, random seed, and temperature are identical to the
original setup, and confirmed that our seed produces deterministic
outputs. CSR detection code is unchanged except for a React patch
adding the \texttt{nextjs-container-errors} class for newer Next.js
versions. The most plausible remaining explanation is regional API
divergence: the original authors used the China Dashscope endpoint
(\texttt{dashscope.aliyuncs.com}), while we use the international
endpoint (\texttt{dashscope-intl.aliyuncs.com}). Alibaba may serve
different model snapshots across regional endpoints under the same model
identifier; we cannot verify this independently. The 72B model is likely robust to
minor weight differences---its code compiles cleanly regardless---while
the 7B model produces borderline outputs where small weight
perturbations can push code from ``barely compiles'' to ``barely
fails,'' amplifying CSR differences disproportionately.

IssAcc shows a similar discrepancy for 7B (paper 0.076 vs.\ ours
0.269). We traced this to the issue extraction logic: if the model's
output does not contain valid JSON inside the \texttt{[ISSUES]} tags,
the parser silently falls back to an empty list, yielding
IssAcc\,=\,0. Minor formatting differences in the 7B output between
endpoints can cause widespread extraction failures, inflating the gap.

At N\,=\,27--28 per framework, even two or three differing samples shift
the mean by several percentage points. We treat the reproduction as
successful based on the close alignment of 72B results, the consistency
of CMLS and CMCS for 7B, and CLIP corroboration across all conditions.


%% ─────────────────────────────────────────
\section{Error Analysis}
% Rubric: systematic analysis, specific examples, failure identification, causal analysis

\subsection{Methodology}

Our proposed approach uses GUI grounding to improve two stages of the
repair pipeline: identifying \textit{what} is visually wrong and
localizing \textit{where} in the code to fix it. We structure our error
analysis around the DesignBench metrics (defined in Section~\ref{sec:baseline})
that correspond to these stages (Table~\ref{tab:metrics}). Issue Accuracy measures whether the model
correctly categorizes the visual defect; CMLS and CMCS measure whether
the model edits the right code in the right way. We additionally report
CSR (whether the output compiles) and CLIP (whether the rendered result
visually resembles the ground truth) as context. Our analysis combines
automated metric computation from DesignBench's evaluation pipeline with
manual inspection of a side-by-side gallery comparing broken input,
ground truth, and model output for each sample. We examine patterns
along two axes: by frontend framework (React, Vue, Angular, Vanilla)
using aggregate metrics, and by defect type (alignment, overflow,
occlusion, crowding, color/contrast) through qualitative comparison.
Sample sizes are small---27--28 per framework and 11--67 per defect
type---and we interpret patterns accordingly.


\subsection{Findings}

The 72B model's results closely match the originally reported values,
with CSR, CMLS, and CMCS all within approximately 5\% on Angular and
visual similarity scores (CLIP) comparable across all frameworks. We
treat this as validation that our reproduction pipeline---prompt
template, rendering environment, and evaluation code---is faithful to
the original benchmark. The remaining analysis focuses on the 7B model,
whose results diverge more substantially and whose failure modes are
directly relevant to our proposed approach.

\paragraph{Compile failures and localization errors compound.}
The 7B model fails to compile 75\% of React outputs and 43\% of Angular
outputs, producing malformed JSX, unterminated strings, and missing
closing tags. However, compilation is not the only bottleneck. Among
samples that do compile, the model frequently edits the wrong code
locations: mean CMLS on compiled samples is 0.28, indicating that even
when the output is syntactically valid, fewer than a third of the
model's AST edits overlap with the ground truth. These are compounding
failures, not independent ones. The model often fails to identify the
visual defect correctly (mean Issue Accuracy of 27\%), and even when
identification succeeds (IssAcc $\geq$ 0.5), CMCS improves only
marginally from 0.354 to 0.374. The correlation between Issue Accuracy
and CMCS is just 0.117, meaning correct defect categorization barely
predicts correct code fixes. The model may recognize that something
looks wrong, attempt a plausible repair strategy, but introduce syntax
errors or edit the wrong element in the process.

\paragraph{Defect type affects difficulty.}
Spatial defects are consistently harder than property-level defects.
Overflow issues have the lowest 7B CMCS (0.220) and the highest compile
failure rate (38.9\%), while occlusion follows closely (CMCS 0.232,
37.9\% compile failures). These defect types require reasoning about
how elements relate to each other in rendered space---which elements
overlap, which exceed their containers---demanding precise understanding
of the CSS box model, flexbox layout, and z-index stacking. In
contrast, color and contrast issues are relatively tractable (CMCS
0.313, 27.3\% compile failures), likely because they involve isolated
property changes rather than spatial reasoning across multiple elements.

\paragraph{Illustrative failure cases.}
We highlight three samples from the side-by-side gallery that
concretize the failure patterns described above.

\begin{figure*}[t]
\centering
\includegraphics[width=\textwidth]{figures/sample11_comparison.png}
\caption{Vanilla/sample\_11 (7B). Left: broken input with low-contrast
sidebar text. Center: ground truth restoring text legibility. Right:
model output, which deletes sidebar content entirely.}
\label{fig:sample11}
\end{figure*}

\textit{Identification failure} (Figure~\ref{fig:sample11}).
In vanilla/sample\_11, the broken UI renders a sidebar with
near-zero contrast---dark text on a dark background. The ground truth
adjusts text color to restore legibility. The 7B model fails to
identify this as a color/contrast defect (IssAcc\,=\,0.00) and instead
removes the sidebar content entirely. Despite achieving perfect
code-level scores (CMLS\,=\,1.0, CMCS\,=\,1.0)---indicating that its
edits target the correct AST nodes---the rendered output is visually
degraded: an empty sidebar rather than a legible one. The model
``repaired'' the contrast problem by eliminating the content that
exhibited it.

\begin{figure*}[t]
\centering
\includegraphics[width=\textwidth]{figures/sample13_comparison.png}
\caption{Vanilla/sample\_13 (72B). Left: broken input with a
multi-column insurance dashboard. Center: ground truth with targeted
layout fixes. Right: model output, which rewrites the page into a
single-column list.}
\label{fig:sample13}
\end{figure*}

\textit{Localization failure} (Figure~\ref{fig:sample13}).
Vanilla/sample\_13 demonstrates the localization bottleneck in the 72B
model. The broken input is an insurance dashboard displaying policy
cards in a multi-column grid. The model partially identifies the defect
(IssAcc\,=\,0.50) but achieves near-zero localization overlap
(CMLS\,=\,0.02): rather than making targeted edits to the existing
layout, it rewrites the page into a single-column list, eliminating the
card grid and the lower purchase section entirely. The model defaults to
regenerating the layout wholesale---precisely the failure mode that GUI
grounding targets, where an explicit mapping from visual regions to code
elements would constrain repairs to the defective components.


\paragraph{Implications for GUI grounding.}
The 7B model is the base architecture for JEDI-7B, the grounding model
we propose to integrate in our approach. The failures documented above
are not fundamental capacity limitations but rather reflect the absence
of structured visual information: the model receives a screenshot and
code side by side but has no explicit mapping between visual regions and
code elements. This is precisely the signal that GUI grounding provides.
Grounding could help the model identify which elements are spatially
misaligned, map those elements to specific code locations, and constrain
repairs to the defective region.

%% ─────────────────────────────────────────
\section{Reflection}

\subsection{Lessons and Capability Gaps}

The reproduction and error analysis surfaced three key lessons about the
task and its evaluation.

First, DesignBench's code-level metrics do not capture visual
correctness. CMLS and CMCS measure Jaccard overlap of AST edit
operations and CodeBLEU of matched edits, respectively. A model that
rewrites a component from scratch but produces a visually identical
result will score near zero on both metrics. This means the benchmark
may systematically undercount functionally correct repairs. We propose
supplementing these metrics with CLIP similarity between rendered
screenshots, which directly measures whether the output \textit{looks}
right regardless of how the code achieves it.

\begin{figure*}[t]
\centering
\includegraphics[width=\textwidth]{figures/sample14_comparison.png}
\caption{Vanilla/sample\_14 (7B). Left: broken input with off-center
profile card text. Center: ground truth with centered text. Right:
model output, which stretches the card to full width without applying
the centering fix.}
\label{fig:sample14}
\end{figure*}

Figure~\ref{fig:sample14} illustrates this disconnect. In
vanilla/sample\_14, the 7B model achieves perfect code-level scores
(IssAcc\,=\,1.0, CMLS\,=\,1.0, CMCS\,=\,1.0) for an off-center
profile card, yet the rendered output stretches the card to full
viewport width rather than centering it. Because CMLS and CMCS only
score edits overlapping the ground-truth set, unmatched edits that
break the layout go unpenalized. CLIP drops to 0.906, capturing the
visual regression that code-level metrics miss and motivating
rendered-screenshot similarity as a supplementary metric.

Second, the 7B model cannot reliably produce syntactically valid code
for framework-based projects (75\% compile failure on React, 43\% on
Angular). This makes it an unsuitable generator for evaluating whether
grounding improves repair quality: if the generator cannot compile its
output, improvements in defect identification are unobservable. Our
error analysis motivates using the 72B model as the generation step,
where compile failures are rare and the analysis can focus on
identification and localization quality.

Third, as shown in Section~\ref{sec:baseline}, defect identification
and code localization are compounding bottlenecks. Correct
identification barely predicts correct fixes. The model lacks an
explicit mapping between visual regions and code elements, which is the
central gap that GUI grounding addresses.
Two strengths of the current baseline deserve attention. The DesignBench
prompt template elicits plausible reasoning from the model---it produces
structured \texttt{[REASONING]} outputs that often correctly describe
the defect even when the resulting fix is wrong. Any approach that
replaces or augments this prompt must preserve this reasoning quality.
Additionally, the single-pass baseline is simple: one model call, no
pipeline orchestration. Our proposed two-step approach (grounding then
generation) introduces complexity and potential error propagation that
must justify itself empirically.

\subsection{Refined Proposal and Alternatives}

Our error analysis refines the approach proposed in our project
proposal. The core idea remains the same: use a GUI grounding model to
produce structured spatial annotations of visual defects, then provide
these annotations to a code-generation model alongside the source code.
The grounding model (JEDI-7B or OmniParser v2) identifies UI elements,
their bounding boxes, and spatial relationships from the rendered
screenshot. The generation model uses these annotations to produce
targeted code repairs.

The key refinement is the choice of generation model. Our original
proposal used Qwen2.5-VL-7B for both grounding and generation. The
error analysis shows this is not viable: the 7B model's compile failures
would mask any improvements from grounding. We now plan to use the 72B
model as the generator, keeping the grounding step at the 7B scale
(where JEDI-7B operates). The comparison becomes: grounding-augmented
identification with 72B generation versus the current single-pass 72B
baseline. Whether identification and generation should be merged into a
single call or separated into distinct steps is itself a variable we
will investigate. We evaluate on the full DesignBench repair task across
all four frontend frameworks.

Our experimental plan proceeds in two stages. First, we evaluate
zero-shot grounding transfer: the grounding model is used as-is, without
task-specific training, to test whether navigation-trained spatial
understanding transfers to defect detection. Second, we apply LoRA
fine-tuning to the grounding model on UI defect data to bridge any
remaining gap between navigation grounding and defect identification. If
grounding shows promise, we plan a further extension using multimodal
RAG to augment the generation step with design guidelines retrieved from
sources such as Material Design, providing normative context that
neither the grounding model nor the generator currently possesses.

If the grounding direction does not improve repair quality, two
alternative approaches are worth exploring. Set-of-marks prompting
overlays numbered visual markers on the screenshot as lightweight
spatial anchors, providing some of the same positional information
without requiring a grounding model. Visual diff guidance takes a
different approach entirely: rendering the broken output, computing
pixel-level differences against a reference, and feeding the diff
regions to the LLM as repair hints. Both are feasible within the project
timeline and attack the visual-to-code gap from different angles.
%% ─────────────────────────────────────────

\bibliographystyle{acl}
\bibliography{latex/references}

\end{document}